\chapter{行列式}
\section{二阶行列式}
略.
\section{三阶行列式}
略.
\section{\texorpdfstring{$n$}{n}阶行列式}
略.
\section{行列式的展开与转置}
4.求证：$n$阶行列式
\[
\begin{vmatrix}
	0 & 0 & \cdots & 0 & b_1 \\
	0 & 0 & \cdots & b_2 & 0 \\
	\vdots & \vdots & &\vdots &\vdots\\
	0 & b_{n- 1} &\cdots &0 & 0\\
	b_n & 0 &\cdots&0&0
\end{vmatrix}
=\left(-1\right)^{\tfrac{1}{2}
n\left(n- 1\right)}
b_1b_2\cdots b_n
\]
\begin{Proof}
对行列式作
相邻列对换
$\displaystyle 
1 +2+ \cdots n -
 1 = \frac{1}{2}n\left(n- 1\right)$
次即证.
\end{Proof}

\section{行列式的计算}
2.计算$n$阶行列式的值
\begin{align*}
	& \text{（1）}
	\begin{vmatrix}
    1 & 1 & \cdots & 1\\
    1 & \binom{2}{1} & \cdots & \binom{n}{1} \\
    1 & \binom{3}{2} & \cdots & \binom{n + 1}{2} \\
    \vdots & \vdots & & \vdots \\
    1 & \binom{n }{n - 1} & \cdots & \binom{2n - 2}{ n - 1}   
\end{vmatrix}
\qquad\qquad\qquad\qquad
\text{（2）}
\begin{vmatrix}
	1 & 2 & 3 & \cdots & n\\
	-1 & 0 & 3 & \cdots & n\\
	-1 & -2 & 0 & \cdots & n\\
	\vdots & \vdots & \vdots & &\vdots\\
	-1 & -2 & -3 & \cdots & 0
\end{vmatrix}
\\
&\text{（3）}
\begin{vmatrix}
	a_1b_1 & a_1b_2 & a_1b_3 & \cdots & a_1b_n\\
	a_1b_2 & a_2b_2 & a_2b_3 & \cdots & a_2b_n\\
	a_1b_3 & a_2b_3 & a_3b_3 & \cdots & a_3b_n\\
	\vdots & \vdots & \vdots &		  & \vdots\\
	a_1b_n & a_2b_n & a_3b_n & \cdots & a_nb_n 
\end{vmatrix}
\qquad\qquad\qquad
\text{（4）}
\begin{vmatrix}
	a & 0 & \cdots & 0 & 1\\
	0 & a & \cdots & 0 & 0\\
	\vdots & \vdots & & \vdots & \vdots\\
	0 & 0 &\cdots &a & 0\\
	1 & 0&\cdots &0&a
\end{vmatrix}
\end{align*}
\begin{solution}
\textup{（1）}因为有
\[
 \binom{n}{m}  - \binom{n - 1}{m - 1} = \binom{n - 1}{m}   
\]
故重复用每一行减去上一行得到
\begin{align*}
	\left|\bm{D}\right|
	=
	\begin{vmatrix}
    1 & 1 & \cdots & 1\\
    1 & \binom{2}{1} & \cdots & \binom{n}{1} \\
    1 & \binom{3}{2} & \cdots & \binom{n + 1}{2} \\
    \vdots & \vdots & & \vdots \\
    1 & \binom{n }{n - 1} & \cdots & \binom{2n - 2}{ n - 1}   
\end{vmatrix}
&= \begin{vmatrix}
	\binom{1}{1} &\binom{2}{1}& \cdots & \binom{n - 1}{1}\\
	\binom{2}{2} & \binom{3}{2}&\cdots & \binom{n}{2}\\
	\vdots &\vdots& &\vdots\\
	\binom{n -1}{n -1}&\binom{n}{n - 1}&\cdots&\binom{2 n  -3}{n - 1}
\end{vmatrix}\\
&=\begin{vmatrix}
	\binom{2}{2} & \cdots & \binom{n-1}{2}\\
	\vdots &&\vdots\\
	\binom{n-1}{n-1}&\cdots&\binom{2n-4}{n-1}
\end{vmatrix}
=\cdots=\begin{vmatrix}
	1&\binom{n-1}{n-2}\\
	1&\binom{n}{n-1}
\end{vmatrix}
=1
\end{align*}
\textup{（2）}
第一行加到其余行即得
\[
\begin{vmatrix}
	1 & 2 & 3 & \cdots & n\\
	-1 & 0 & 3 & \cdots & n\\
	-1 & -2 & 0 & \cdots & n\\
	\vdots & \vdots & \vdots & &\vdots\\
	-1 & -2 & -3 & \cdots & 0
\end{vmatrix}
=
\begin{vmatrix}
	1 & 2 & 3 & \cdots & n\\
	0 & 2 & 6 & \cdots & 2n\\
	0 & 0 & 3 & \cdots & 2n\\
	\vdots & \vdots & \vdots & &\vdots\\
	0 & 0 & 0 & \cdots & n
\end{vmatrix}
=n!
\]
\textup{（3）}
简单的列变换
\begin{align*}
	\begin{vmatrix}
	a_1b_1 & a_1b_2 & a_1b_3 & \cdots & a_1b_n\\
	a_1b_2 & a_2b_2 & a_2b_3 & \cdots & a_2b_n\\
	a_1b_3 & a_2b_3 & a_3b_3 & \cdots & a_3b_n\\
	\vdots & \vdots & \vdots &		  & \vdots\\
	a_1b_n & a_2b_n & a_3b_n & \cdots & a_nb_n 
\end{vmatrix}
&=a_1b_n
\begin{vmatrix}
	b_1 & a_1b_2 & a_1b_3 & \cdots & a_1b_n\\
	b_2 & a_2b_2 & a_2b_3 & \cdots & a_2b_n\\
	b_3 & a_2b_3 & a_3b_3 & \cdots & a_3b_n\\
	\vdots & \vdots & \vdots &		  & \vdots\\
	1 & a_2 & a_3 & \cdots & a_n 
\end{vmatrix}\\
&=a_1b_n
\begin{vmatrix}
	b_1 & a_1b_2 - a_2b_1 & a_1b_3 & \cdots & a_1b_n\\
	b_2 & 0 & a_2b_3 & \cdots & a_2b_n\\
	b_3 & 0 & a_3b_3 & \cdots & a_3b_n\\
	\vdots & \vdots & \vdots &		  & \vdots\\
	1 & 0 & a_3 & \cdots & a_n 
\end{vmatrix}\\
&=a_1b_n
\begin{vmatrix}
	b_1 & a_1b_2 - a_2b_1 & a_1b_3-a_3b_1 & \cdots & a_1b_n-a_nb_1\\
	b_2 & 0 & a_2b_3-a_3b_2 & \cdots & a_2b_n-a_nb_2\\
	b_3 & 0 & 0 & \cdots & a_3b_n-a_nb_3\\
	\vdots & \vdots & \vdots &		  & \vdots\\
	b_{n-1}&0&0&\cdots&a_{n-1}b_n-a_nb_{n-1}\\
	1 & 0 & 0 & \cdots & 0 
\end{vmatrix}
\\
&=\left(-1\right)^{n+1}a_1b_n
\begin{vmatrix}
	a_1b_2 - a_2b_1 & a_1b_3-a_3b_1 & \cdots & a_1b_n-a_nb_1\\
	 0 & a_2b_3-a_3b_2 & \cdots & a_2b_n-a_nb_2\\
	 0 & 0 & \cdots & a_3b_n-a_nb_3\\
	 \vdots & \vdots &		  & \vdots\\
	0&0&\cdots&a_{n-1}b_n-a_nb_{n-1}
\end{vmatrix}\\
&=\left(-1\right)^{n+1}a_1b_n
\prod _{i = 1}^{n-1}\left(
	a_ib_{i+1}-a_{i+1}b_i
\right)\\
&=a_1b_n
\prod _{i = 1}^{n-1}\left(
	a_{i+1}b_{i}-a_{i}b_{i+1}
\right)
\end{align*}
\textup{（4）}
拆分即可
\begin{align*}
	\begin{vmatrix}
	a & 0 & \cdots & 0 & 1\\
	0 & a & \cdots & 0 & 0\\
	\vdots & \vdots & & \vdots & \vdots\\
	0 & 0 &\cdots &a & 0\\
	1 & 0&\cdots &0&a
\end{vmatrix}
&=\begin{vmatrix}
	a & 0 & \cdots & 0 & 0\\
	0 & a & \cdots & 0 & 0\\
	\vdots & \vdots & & \vdots & \vdots\\
	0 & 0 &\cdots &a & 0\\
	1 & 0&\cdots &0&a
\end{vmatrix}
+
\begin{vmatrix}
	0 & 0 & \cdots & 0 & 1\\
	0 & a & \cdots & 0 & 0\\
	\vdots & \vdots & & \vdots & \vdots\\
	0 & 0 &\cdots &a & 0\\
	1 & 0&\cdots &0&a
\end{vmatrix}\\
&=a^n+\left(-1\right)^{n+1}
\left(-1\right)^na^{n-2}
=a^n-a^{n-2}
\end{align*}
\end{solution}

3.设$b_{ij}=\left(a_{i1}+\cdots
+a_{in}\right)-a_{ij}$,求证
\[
\begin{vmatrix}
	b_{11} & \cdots & b_{1n}\\
	\vdots &&\vdots\\
	b_{n1}&\cdots &b_{nn}
\end{vmatrix}
=\left(-1\right)^{n-1}\left(n-1\right)
\begin{vmatrix}
	a_{11} & \cdots & a_{1n}\\
	\vdots &&\vdots\\
	a_{n1}&\cdots &a_{nn}
\end{vmatrix}
\]
\begin{Proof}
	容易有
	\begin{align*}
		\begin{vmatrix}
			b_{11} & \cdots & b_{1n}\\
	\vdots &&\vdots\\
	b_{n1}&\cdots &b_{nn}
		\end{vmatrix}=
		\begin{vmatrix}
			\left(a_{11}+\cdots
			+a_{1n}\right)-a_{11}
			&\cdots&\left(a_{11}+
			\cdots
			+a_{1n}\right)-a_{1n}\\
			\vdots&&\vdots\\
			\left(a_{n1}+\cdots
			+a_{nn}\right)-a_{n1}
			&\cdots&\left(a_{n1}+
			\cdots
			+a_{nn}\right)-a_{nn}
		\end{vmatrix}
	\end{align*}
	将其余列加到第一列
	\begin{align*}
		&\begin{vmatrix}
			\left(a_{11}+\cdots
			+a_{1n}\right)-a_{11}
			&\cdots&\left(a_{11}+
			\cdots
			+a_{1n}\right)-a_{1n}\\
			\vdots&&\vdots\\
			\left(a_{n1}+\cdots
			+a_{nn}\right)-a_{n1}
			&\cdots&\left(a_{n1}+
			\cdots
			+a_{nn}\right)-a_{nn}
		\end{vmatrix}\\
		=&
		\begin{vmatrix}
			\left(n-1\right)\left(a_{11}+\cdots
			+a_{1n}\right)
			&\cdots&\left(a_{11}+
			\cdots
			+a_{1n}\right)-a_{1n}\\
			\vdots&&\vdots\\
			\left(n-1\right)\left(a_{n1}+\cdots
			+a_{nn}\right)
			&\cdots&\left(a_{n1}+
			\cdots
			+a_{nn}\right)-a_{nn}
		\end{vmatrix}
		\\
		=&\left(n-1\right)
		\begin{vmatrix}
			\left(a_{11}+\cdots
			+a_{1n}\right)
			&\cdots&\left(a_{11}+
			\cdots
			+a_{1n}\right)-a_{1n}\\
			\vdots&&\vdots\\
			\left(a_{n1}+\cdots
			+a_{nn}\right)
			&\cdots&\left(a_{n1}+
			\cdots
			+a_{nn}\right)-a_{nn}
		\end{vmatrix}
	\end{align*}
	再用后几列分别减去第一列,然后其余列加到
	第一列
	\begin{align*}
		&\left(n-1\right)
		\begin{vmatrix}
			\left(a_{11}+\cdots
			+a_{1n}\right)
			&\cdots&\left(a_{11}+
			\cdots
			+a_{1n}\right)-a_{1n}\\
			\vdots&&\vdots\\
			\left(a_{n1}+\cdots
			+a_{nn}\right)
			&\cdots&\left(a_{n1}+
			\cdots
			+a_{nn}\right)-a_{nn}
		\end{vmatrix}
		\\=&\left(n-1\right)
		\begin{vmatrix}
			a_{11} & \cdots & -a_{1n}\\
	\vdots &&\vdots\\
	a_{n1}&\cdots &-a_{nn}
		\end{vmatrix}=
		\left(-1\right)^{n-1}\left(n-1\right)
\begin{vmatrix}
	a_{11} & \cdots & a_{1n}\\
	\vdots &&\vdots\\
	a_{n1}&\cdots &a_{nn}
\end{vmatrix}
\qedhere
	\end{align*}
\end{Proof}
4.利用Vandermonde行列式求
\[
\begin{vmatrix}
	a_1^{n-1}&a_1^{n-2}b_1&\cdots&a_1b_1^{n-2}&b_1^{n-1}\\
	a_2^{n-1}&a_2^{n-2}b_2&\cdots&a_2b_2^{n-2}&b_2^{n-1}\\
	\vdots&\vdots&&\vdots&\vdots\\
	a_n^{n-1}&a_n^{n-2}b_n&\cdots&a_nb_n^{n-2}&b_n^{n-1}
\end{vmatrix}
\]
\begin{solution}
	明显地,若$a_i\neq 0\left(
		1\leqslant 
		i \leqslant n
	\right)$
	\begin{align*}
		&\begin{vmatrix}
	a_1^{n-1}&a_1^{n-2}b_1&\cdots&a_1b_1^{n-2}&b_1^{n-1}\\
	a_2^{n-1}&a_2^{n-2}b_2&\cdots&a_2b_2^{n-2}&b_2^{n-1}\\
	\vdots&\vdots&&\vdots&\vdots\\
	a_n^{n-1}&a_n^{n-2}b_n&\cdots&a_nb_n^{n-2}&b_n^{n-1}
	\end{vmatrix}\\=
&\left(a_1\cdots a_n\right)
^{n-1}
	\begin{vmatrix}
	1&\cfrac{b_1}{a_1}&\cdots&\left(
		\cfrac{b_1}{a_1}
	\right)^{n-2}&\left(
		\cfrac{b_1}{a_1}
	\right)^{n-1}\\
	1&\cfrac{b_2}{a_2}&\cdots&\left(
		\cfrac{b_2}{a_2}
	\right)^{n-2}&\left(
		\cfrac{b_2}{a_2}
	\right)^{n-1}\\
	\vdots&\vdots&&\vdots&\vdots\\
	1&\cfrac{b_n}{a_n}&\cdots&\left(
		\cfrac{b_n}{a_n}
	\right)^{n-2}&\left(
		\cfrac{b_n}{a_n}
	\right)^{n-1}
	\end{vmatrix}\\
	=&\left(a_1\cdots a_n\right)
	^{n-1}
	\prod _{1\leqslant i
	< j \leqslant n}\left(
		\frac{b_j}{a_j}-
		\frac{b_i}{a_i}
	\right)\\
	=&\prod _{1\leqslant i
	< j \leqslant n}
	\left(a_ib_j-a_jb_i\right)
\end{align*}
若只有一个$a_i=0$,按第$i$行展开
得到的行列式具有相同类型；
若至少有两个$a_i=a_j=0$,则行列式为零,
容易证明,后两种情况是与第一
种情况统一的.
\end{solution}

5.设$f_i\left(x\right)\left(
	i=1,2,\cdots,n
\right)$均为次数不超过$n-2$的关于$x$的多项式,求证：
对任意$n$个数$a_1,a_2,\cdots,a_n$,均有
\[
\begin{vmatrix}
	f_1\left(a_1\right) & 
	f_1\left(a_2\right) & \cdots &
	f_1\left(a_n\right)\\
	f_2\left(a_1\right) & 
	f_2\left(a_2\right) & \cdots &
	f_2\left(a_n\right)\\
	\vdots&\vdots&&\vdots\\
	f_n\left(a_1\right) & 
	f_n\left(a_2\right) & \cdots &
	f_n\left(a_n\right)\\
\end{vmatrix}
=0
\]
提示：试证明
\[
\begin{vmatrix}
	f_1\left(x\right) & 
	f_1\left(a_2\right) & \cdots &
	f_1\left(a_n\right)\\
	f_2\left(x\right) & 
	f_2\left(a_2\right) & \cdots &
	f_2\left(a_n\right)\\
	\vdots&\vdots&&\vdots\\
	f_n\left(x\right) & 
	f_n\left(a_2\right) & \cdots &
	f_n\left(a_n\right)\\
\end{vmatrix}
\equiv 0
\]
\begin{Proof}
	构造
	\begin{align*}
		g\left(x\right)
		=\begin{vmatrix}
			f_1\left(x\right) & 
			f_2\left(x\right) & \cdots &
			f_n\left(x\right)\\
			f_1\left(a_2\right) & 
			f_2\left(a_2\right) & \cdots &
			f_n\left(a_2\right)\\
			\vdots&\vdots&&\vdots\\
			f_1\left(a_n\right) & 
			f_2\left(a_n\right) & \cdots &
			f_n\left(a_n\right)\\
		\end{vmatrix}
	\end{align*}
	容易有
	$a_i=a_j\left(2\leqslant i\neq j\leqslant 
	n\right)$时$g\left(x\right)\equiv 0$
	；$a_i\neq a_j\left(2\leqslant i\neq j\leqslant 
	n\right)$时$g\left(a_i\right)
	=0\left(2\leqslant i
	\leqslant n\right)$,又$\mathrm{deg}f_i\left(
		x
	\right)\leqslant n-2$,故
	$g\left(x\right)\equiv 0$.得证.
	\qedhere
\end{Proof}
另一种更自然的证法如下
\begin{Proof}
	因为$f_k\left(x\right)$均为次数不超过
	$n-2$的多项式,
	所以按照单项式
	一直做行列式的
	拆分一定可以使得
	原行列式等于若干个元素均为
	单项式的行列式之和,根据抽屉原理,
	每个行列式中
	至少有两列元素共用同一个单项式,
	于是这些行列式值均为零,得证.
	\qedhere
\end{Proof}

6.设$t$为参数,
\[
\left|\bm{A}\left(t\right)\right|
=\begin{vmatrix}
	a_{11}+t & a_{12}+t & \cdots &a_{1n}+t\\
	a_{21}+t & a_{22}+t & \cdots &a_{2n}+t\\
	\vdots &\vdots&&\vdots
	\\
	a_{n1}+t & a_{n2}+t & \cdots &a_{nn}+t
\end{vmatrix}
\]
求证
\[
\left|\bm{A}\left(t\right)\right|
=\left|\bm{A}\left(0\right)\right|
+t\sum_{i,j=1}^{n}A_{ij}
\]
其中$A_{ij}$为$a_{ij}$在
$\left|\bm{A}\left(0\right)\right|$中的
代数余子式.
\begin{Proof}
	我们证明更强的结论,设
	\[
	\left|\bm{A}\right|
	=\begin{vmatrix}
	a_{11} & \cdots & a_{1n}\\
	\vdots &\ddots&\vdots\\
	a_{n1}&\cdots &a_{nn}
\end{vmatrix}
	\]
	则
	\[
	\left|\bm{A}\left(
		t_1,t_2,\cdots,t_n
	\right)\right|
	=\begin{vmatrix}
	a_{11}+t_1 & a_{12}+t_2&
	\cdots & a_{1n}+t_n\\
	a_{21}+t_1 & a_{22}+t_2&
	\cdots & a_{2n}+t_n\\
	\vdots &\vdots&&\vdots\\
	a_{n1}+t_1&a_{n2}+t_2&
	\cdots &a_{nn}+t_n
\end{vmatrix}
=\left|\bm{A}\right|
+\sum_{j=1}^{n}\left(
	t_j\sum_{i=1}^{n}
	A_{ij}
\right)
	\]
	将上式第一列拆开
	\begin{align*}
		\left|\bm{A}\left(
		t_1,t_2,\cdots,t_n
	\right)\right|
	&=\begin{vmatrix}
	a_{11}+t_1 & a_{12}+t_2&
	\cdots & a_{1n}+t_n\\
	a_{21}+t_1 & a_{22}+t_2&
	\cdots & a_{2n}+t_n\\
	\vdots &\vdots&&\vdots\\
	a_{n1}+t_1&a_{n2}+t_2&
	\cdots &a_{nn}+t_n
\end{vmatrix}\\
&=\begin{vmatrix}
	a_{11} & a_{12}+t_2&
	\cdots & a_{1n}+t_n\\
	a_{21} & a_{22}+t_2&
	\cdots & a_{2n}+t_n\\
	\vdots &\vdots&&\vdots\\
	a_{n1}&a_{n2}+t_2&
	\cdots &a_{nn}+t_n
\end{vmatrix}+
\begin{vmatrix}
	t_1 & a_{12}+t_2&
	\cdots & a_{1n}+t_n\\
	t_1 & a_{22}+t_2&
	\cdots & a_{2n}+t_n\\
	\vdots &\vdots&&\vdots\\
	t_1&a_{n2}+t_2&
	\cdots &a_{nn}+t_n
\end{vmatrix}\\
&=\cdots=
\left|\bm{A}\right|
+\sum_{j=1}^{n}\left(
	t_j\sum_{i=1}^{n}
	A_{ij}
	\right)
	\end{align*}
如此得证.
	\qedhere
\end{Proof}
\section{行列式的等价定义}
4.若一个$n$阶行列式中零元素个数超过
$n^2-n$,证明：该行列式值为零.
\begin{Proof}
	由组合定义
	\[
	\left|\bm{A}\right|=
	\sum_{\left(k_1,k_2,\cdots,k_n\right)\in S_n}
	\left(-1\right)^{N\left(k_1,k_2,\cdots,k_n\right)}a_{1k_1}a_{2k_2}\cdots
	a_{nk_n}
	\]
	由于行列式$\left|\bm{A}\right|$中零元素的个数超过
$n^2-n$个,故$a_{1k_1},a_{2k_2},\cdots,
	a_{nk_n}$中至少一个为零,于是$a_{1k_1}a_{2k_2}\cdots
	a_{nk_n}=0$,故$\left|\bm{A}\right|=0.$或者考虑抽屉原理,至少有
	一列其零元素个数不少于
	\[
	\left[\frac{n^2-n}{n}\right]+1=n
	\]
	于是至少有一列元素全为零,故$\left|\bm{A}\right|=0.$
\end{Proof}

5.设$f_{ij}\left(x\right)$是可导函数,记
\[
F\left(x\right)=
\begin{vmatrix}
    f_{11}\left(x\right)&f_{12}\left(x\right)&\cdots&f_{1n}\left(x\right)\\
    f_{21}\left(x\right)&f_{22}\left(x\right)&\cdots&f_{2n}\left(x\right)\\
    \vdots&\vdots&&\vdots\\
    f_{n1}\left(x\right)&f_{n2}\left(x\right)&\cdots&f_{nn}\left(x\right)
\end{vmatrix}
\]
证明：
\[
\frac{\mathrm{d}}{\mathrm{d}x}F\left(x\right)
=
\sum_{i=1}^{n}F_i\left(x\right)
\]
其中
\[
F_i\left(x\right)
=
\begin{vmatrix}
    f_{11}\left(x\right)&f_{12}\left(x\right)&\cdots&\cfrac{\mathrm{d}}{\mathrm{d}x}f_{1i}\left(x\right)&\cdots&f_{1n}\left(x\right)\\
    f_{21}\left(x\right)&f_{22}\left(x\right)&\cdots&\cfrac{\mathrm{d}}{\mathrm{d}x}f_{2i}\left(x\right)&\cdots&f_{2n}\left(x\right)\\
    \vdots&\vdots&&\vdots&&\vdots\\
    f_{n1}\left(x\right)&f_{n2}\left(x\right)&\cdots&\cfrac{\mathrm{d}}{\mathrm{d}x}f_{ni}\left(x\right)&\cdots&f_{nn}\left(x\right)
\end{vmatrix}
\]
\begin{Proof}
考虑数学归纳法
,对阶数$n$归纳,$n=1$时显然成立.
设结论对$n-1$阶行列式成立
,将行列式展开
\[
F\left(x\right)
=
f_{11}\left(x\right)
A_{11}\left(x\right)
+f_{21}\left(x\right)A_{21}\left(x\right)+
\cdots+f_{n1}\left(x\right)A_{n1}\left(x\right)
\]
其中$A_{i1}$是$f_{i1}\left(x\right)$的代数余子式,
两端求导并记$A_{ij}^k\left(x\right)$
是对$A_{ij}\left(x\right)$第$k$列元素求导后得到的行列式
,则
\begin{align*}
    \frac{\mathrm{d}}{\mathrm{d}x}F\left(x\right)
    &=\frac{\mathrm{d}}{\mathrm{d}x}\left(
        \sum_{i=1}^{n}f_{i1}\left(x\right)A_{i1}\left(x\right)
    \right)
    =\sum_{i=1}^{n}f'_{i1}\left(x\right)A_{i1}\left(x\right)
    +\sum_{i=1}^{n}f_{i1}\left(x\right)A'_{i1}\left(x\right)\\
    &=F_1\left(x\right)+\sum_{i=1}^{n}f_{i1}\left(x\right)
    \left(\sum_{k=1}^{n-1}A^k_{i1}\left(x
    \right)\right)=F_1\left(x\right)+
    \sum_{k=1}^{n-1}\sum_{i=1}^{n}
    f_{i1}\left(x\right)A^k_{i1}\left(x\right)
\end{align*}
又$\displaystyle
\sum_{i=1}^{n}f_{i1}\left(x\right)A_{i1}^k\left(x
\right)=F_{k+1}\left(x\right)\left(
    1\leqslant k\leqslant n-1
\right)$,故
\[
\frac{\mathrm{d}}{\mathrm{d}x}F\left(x\right)
=
F_{1}\left(x\right)+F_2\left(x
\right)+\cdots+F_n\left(x\right)
\qedhere
\]
\end{Proof}
利用组合定义也可简单地证明.
\[
F\left(x\right)=
\sum_{\left(k_1,k_2,\cdots,k_n\right)\in
S_n}\left(-1\right)^{N\left(k_1,k_2,\cdots,k_n\right)}
f_{k_11}\left(x\right)f_{k_22}\left(x\right)\cdots
f_{k_nn}\left(x\right).
\]
此处不再赘述.

6.设
\[
f\left(x\right)=\begin{vmatrix}
	x-a_{11}&-a_{12}&\cdots&-a_{1n}\\
	-a_{21}&x-a_{22}&\cdots&-a_{2n}\\
	\vdots&\vdots&&\vdots\\
	-a_{n1}&-a_{n2}&\cdots&x-a_{nn}
\end{vmatrix}
\]其中$x$为未定元,$a_{ij}$为常数.证明：
$f\left(x\right)$是一个最高次项系数为1的$n$次多项式
,且其$n-1$次项的系数为$-\left(a_{11}+a_{22}+\cdots+a_{nn}\right).$
\begin{Proof}
由组合定义,$f\left(x\right)$的最高次项出现于$\left(x-a_{11}\right)\left(x-a_{22}
\right)\cdots\left(x-a_{nn}\right)$单项,且其他单项对整体仅贡献不高于$n-2$次的项,
于是$f\left(x\right)$是一个最高次项系数为1的$n$次多项式
,且其$n-1$次项的系数为$-\left(a_{11}+a_{22}+\cdots+a_{nn}\right).$
\end{Proof}
\section{Laplace定理}
2.设$\omega$为三次单位根
\[
\omega = - \frac{1}{2}+\frac{\sqrt{3}}{2}\mathrm{i}
\]
求
\[
\begin{vmatrix}
	1&1&0&0&1&0\\
	1&\omega&0&0&\omega^2&0\\
	a_1&b_1&1&1&c_1&1\\
	a_2&b_2&1&\omega^2&c_2&\omega\\
	a_3&b_3&1&\omega&c_3&\omega^2\\
	1&\omega^2&0&0&\omega&0
\end{vmatrix}
\]
\begin{solution}
	取定\textup{1、2、6}行列,注意到只有$j_1=1,j_2=2,j_3=5$的子式
	不为零,于是
	\begin{align*}
		\begin{vmatrix}
	1&1&0&0&1&0\\
	1&\omega&0&0&\omega^2&0\\
	a_1&b_1&1&1&c_1&1\\
	a_2&b_2&1&\omega^2&c_2&\omega\\
	a_3&b_3&1&\omega&c_3&\omega^2\\
	1&\omega^2&0&0&\omega&0
\end{vmatrix}=-\begin{vmatrix}
	1&1&1\\
	1&\omega&\omega^2\\
	1&\omega^2&\omega
\end{vmatrix}\begin{vmatrix}
	1&1&1\\
	1&\omega^2&\omega\\
	1&\omega&\omega^2
\end{vmatrix}
	\end{align*}
	又因为
	\[
	\begin{cases*}
	\omega^3=1\\
	\omega^2+\omega+1=0
	\end{cases*}
	\]
	故
	\begin{align*}
		-\begin{vmatrix}
	1&1&1\\
	1&\omega&\omega^2\\
	1&\omega^2&\omega
\end{vmatrix}\begin{vmatrix}
	1&1&1\\
	1&\omega^2&\omega\\
	1&\omega&\omega^2
\end{vmatrix}&=-9\begin{vmatrix}
	1&0&0\\
	1&\omega&\omega^2\\
	1&\omega^2&\omega
\end{vmatrix}\begin{vmatrix}
	1&0&0\\
	1&\omega^2&\omega\\
	1&\omega&\omega^2
\end{vmatrix}\\
&=-9\left(\omega^2-\omega\right)\left(
	\omega-\omega^2
\right)\\
&=9\left(\omega^2+\omega-2\right)
\\
&=-27
	\end{align*}
\end{solution}

3.证明：
	确定代数余子式的符号时,无需利用该子式的行列号码和,
	利用余子式行列号码和即可.
	\begin{Proof}
		两组号码和之和为$2\left(1+2+\cdots+n\right)$为偶数.
	\end{Proof}

4.证明：
\[
\left(ab'-ba'\right)\left(cd'-dc'\right)-\left(
	ac'-ca'
\right)\left(bd'-db'\right)+\left(
	ad'-da'
\right)\left(bc'-cb'\right)=0
\]
\begin{Proof}
	考虑行列式
\[
\begin{vmatrix}
	a&a'&a&a'\\
	b&b'&b&b'\\
	c&c'&c&c'\\
	d&d'&d&d'
\end{vmatrix}=0
\]
取定第一、二列展开得证.
\end{Proof}

5.求$2n$阶行列式的值
\[
\begin{vmatrix}
	a&&&&&b\\
	&\ddots&&&\iddots &\\
	&&a&b&&\\
	&&b&a&&\\
	&\iddots &&&\ddots &\\
	b&&&&&a
\end{vmatrix}
\]
\begin{solution}
	对第一行及最后一行不断用\textup{Laplace}定理：
	明显每次只能取第一行、第一列、最后一行、最后一列,
	且号码和均为偶数,那么结果为
	\[
	\left(a^2-b^2\right)^n
	\]
\end{solution}
\section{复习题一}
7.求证：$n$阶行列式\[\left|\bm{A}\right|=
\begin{vmatrix}
    \cos x&1&0&0&\cdots&0&0&0\\
    1&2\cos x&1&0&\cdots&0&0&0\\
    0&1&2\cos x&1&\cdots&0&0&0\\
    \vdots&\vdots&\vdots&\vdots&&\vdots&\vdots&\vdots\\
    0&0&0&0&\cdots&1&2\cos x&1\\
    0&0&0&0&\cdots&0&1&2\cos x
\end{vmatrix}=\cos nx
\]
\begin{Proof}
    因为\begin{align*}
        \left|\bm{A}\right|&=\begin{vmatrix}
    2\cos x&1&0&0&\cdots&0&0&0\\
    1&2\cos x&1&0&\cdots&0&0&0\\
    0&1&2\cos x&1&\cdots&0&0&0\\
    \vdots&\vdots&\vdots&\vdots&&\vdots&\vdots&\vdots\\
    0&0&0&0&\cdots&1&2\cos x&1\\
    0&0&0&0&\cdots&0&1&2\cos x
\end{vmatrix}\\
&-\begin{vmatrix}
    \cos x&1&0&0&\cdots&0&0&0\\
    0&2\cos x&1&0&\cdots&0&0&0\\
    0&1&2\cos x&1&\cdots&0&0&0\\
    \vdots&\vdots&\vdots&\vdots&&\vdots&\vdots&\vdots\\
    0&0&0&0&\cdots&1&2\cos x&1\\
    0&0&0&0&\cdots&0&1&2\cos x
\end{vmatrix}\\
&=D_n-\cos xD_{n-1}
    \end{align*}其中$D_n$是$a_1=a_2=\cdots=a_n=2\cos x,b_1=b_2=\cdots=b_{n-1}=c_1=c_2=\cdots=c_{n-1}=1$的三对角行列式,于是\[
    D_n=2\cos xD_{n-1}-D_{n-2},D_0=1,D_1=2\cos x
    \]即\[
    \frac{D_n-\alpha D_{n-1}}{D_{n-1}-\alpha D_{n-2}}=\beta
    \]其中$\alpha=\cos x+\mathrm{i}\sin x,\beta=\cos x-\mathrm{i}\sin x.$于是$x\neq k\pi,k\in\mathbb{Z}\Longrightarrow\alpha\neq\beta$时\begin{align*}
        D_n&=\frac{
            \alpha^{n+1}-\beta^{n+1}
        }{
            \alpha-\beta
        }\\
        &=\frac{
            \sin\left(n+1\right)x
        }{\sin x}
    \end{align*}则\begin{align*}
        \left|\bm{A}\right|&=D_n-\cos xD_{n-1}\\
        &=\cos nx
    \end{align*}而$x=k\pi,k\in\mathbb{Z}\Longrightarrow\alpha=\beta$时\begin{align*}
        D_n=\left(n+1\right)\cos^{n}x=\left(n+1\right)\left(-1\right)^{kn},k\in\mathbb{Z}
    \end{align*}于是\begin{align*}
        \left|\bm{A}\right|&=D_n-\cos xD_{n-1}\\
        &=\left(n+1\right)\left(-1\right)^{kn}-\left(-1\right)^kn\left(-1\right)^{k\left(n-1\right)}\\
        &=\left(-1\right)^{kn}=\cos nx\qedhere
    \end{align*}
\end{Proof}
8.计算行列式\[
\left|\bm{A}\right|=\begin{vmatrix}
    x_1&y&y&\cdots&y&y\\
    z&x_2&y&\cdots&y&y\\
    z&z&x_3&\cdots&y&y\\
    \vdots&\vdots&\vdots&&\vdots&\vdots\\
    z&z&z&\cdots&x_{n-1}&y\\
    z&z&z&\cdots&z&x_n
\end{vmatrix}
\]\begin{solution}
    因为\begin{align*}
        \left|\bm{A}\left(t\right)\right|&=\begin{vmatrix}
            a_{11}+t&a_{12}+t&a_{13}+t&\cdots&a_{1n}+t\\
            a_{21}+t&a_{22}+t&a_{23}+t&\cdots&a_{2n}+t\\
            \vdots&\vdots&\vdots&&\vdots\\
            a_{n1}+t&a_{n2}+t&a_{n3}+t&\cdots&a_{nn}+t
        \end{vmatrix}\\
        &=\left|\bm{A}\left(0\right)\right|+t\sum_{i, j=1}^{n}A_{ij}
    \end{align*}则由题有 $y\neq z$时\[
    \left|\bm{A}\right|=\frac{1}{y-z}\left[
        y\prod_{i=1}^{n}\left(x_i-z\right)-z\prod_{i=1}^{n}\left(x_i-y\right)
    \right]
    \]$y=z$时由递推\[
    D_n=\left(x_n-y\right)D_{n-1}+y\prod_{i=1}^{n-1}\left(x_i-y\right)
    \]得到\[
    D_n=\prod_{i=1}^{n}\left(x_i-y\right)+y\sum_{i=1}^{n}\prod_{i\neq j}\left(x_j-y\right)
    \]
\end{solution}